\section{Related work}%
\label{sec:related}

Novel view synthesis (NVS) is a challenging and ill-posed problem that can be categorized into two main categories: view interpolation methods~\cite{niemeyer2022regnerf, chen2021mvsnerf, chen_single-stage_2023, chen2024mvsplat, chen2024mvsplat360, wu2024reconfusion, truong2023sparf, gao2024cat3d}, which generate views close to given inputs, and view extrapolation or autoregressive view generation, where novel views extend significantly beyond the original scene, introducing substantial new content.
The latter is particularly difficult from single images.
Between these lie single-view scene reconstruction methods, which succeed mainly for single-object scenes~\cite{melas-kyriazi23realfusion, shi23zero123:, hong24lrm:, szymanowicz24splatter, boss2024sf3d, xiang2024structured} or highly bounded scenes~\cite{szymanowicz2024flash3d}.
However, their extrapolation capabilities remain limited.
Most relevant to our work are single-image view extrapolation models, which fall into two categories: those incorporating explicit geometric modeling with inpainting, and those based directly on image or video generation.

\paragraph{Inpainting-based view extrapolation.}
Inpainting-based methods~\cite{wiles2020synsin, SceneScape, hoellein2023text2room, muller2024multidiff, popov24camctrl3d, yu2024viewcrafter, seo2024genwarp, koh2023simple, liu2021infinite, rockwell2021pixelsynth, yu23wonderjourney:, yu25wonderworld:} use pre-trained 3D reconstruction models to generate 3D representations—such as meshes, point clouds, or Gaussian splats—from images.
These representations are reprojected into novel views, where 2D inpainting fills missing regions.
SceneScape~\cite{SceneScape} reconstructs meshes from images using pre-trained depth estimators.
Diffusion-based inpainting then completes the projected novel view, which is reprojected back to refine and extend the mesh.
This process is iterated to generate novel views.
MultiDiff~\cite{muller2024multidiff} uses depth estimators to warp reference images into multiple novel views, training diffusion models to inpaint missing regions simultaneously.
CamCtrl3D~\cite{popov24camctrl3d} follows a similar approach but adds ray maps to condition diffusion models for view synthesis.
ViewCrafter~\cite{yu2024viewcrafter} uses pre-trained pointmap estimators to create point clouds and trains video diffusion models for inpainting.
GenWarp~\cite{seo2024genwarp} warps 2D coordinate representations of reference views into novel views and uses them to condition generation.

While these methods improve geometric consistency, they are susceptible to errors in depth or pointmap estimation, which propagate distortions across generated views.
Once 3D representations are constructed, inaccuracies from depth or pointmap estimation become difficult to correct.
Moreover, scaling to large scenes is computationally intensive, as storing and processing high-fidelity 3D representations requires substantial memory.

\begin{figure*}[h]
\centering
\includegraphics[width=0.95\textwidth,clip,trim=0 0.4cm 0 0cm]{imgs/surfel_demo.pdf}
\caption{\textbf{Surfel-based memory index.}
Each surfel stores indices of views that observed it.
We color-code each surfel by contributing view indices.
This spatial index enables retrieval of relevant past views: when generating a novel view, we identify visible surfels from the target viewpoint and retrieve views that previously observed those same regions, naturally accounting for occlusion.
}%
\label{fig:surfel_demo}
\end{figure*}


\paragraph{Multi-view-based view extrapolation.}
Another category~\cite{geogpt, ren2022look, tseng2023consistent, photometricnvs, wang2024motionctrl, gao2024cat3d, seva} avoids explicitly modeling 3D scene geometry.
Instead, these approaches generate novel views by conditioning on previously rendered or generated views.
GeoGPT~\cite{geogpt} uses autoregressive likelihood models to synthesize novel views from single images.
Building on this, LookOut~\cite{ren2022look} improves consistency by generating view sequences along predefined trajectories, conditioning on up to two preceding views.
Other methods~\cite{tseng2023consistent, photometricnvs} enhance novel view synthesis by employing cross-view attention or enforcing epipolar constraints within diffusion models.
MotionCtrl~\cite{wang2024motionctrl} leverages geometric priors from pre-trained video diffusion models to synthesize camera-controllable videos.
These methods are \textit{trajectory-based}, requiring input views to follow continuous spatial paths.
In contrast, \textit{image-set models} like CAT3D~\cite{gao2024cat3d} and more recently SEVA~\cite{seva} synthesize novel views from sparse and unordered input images, without trajectory constraints.
However, these methods are computationally expensive, with $O(n^2)$ complexity for attention, limiting conditioning views to small context windows containing only recently generated frames.
This limitation impacts long-term consistency of generated images.

Several concurrent works explore memory-based approaches for scene generation.
Genie~\cite{bruce2024genie, parker-holder24genie} maintains memory through recurrent state features.
Gen3C~\cite{ren2025gen3c} conditions on stored point clouds like ViewCrafter~\cite{yu2024viewcrafter}, inheriting inpainting-based limitations.
StarGen~\cite{zhai2025stargen} and WorldMem~\cite{xiao2025worldmemlongtermconsistentworld} use distance- and field-of-view-based spatial memory retrieval, which can recover spatially correlated views but lack geometric reasoning for occlusions.
In contrast, our surfel-based memory explicitly models occlusions and provides more principled view selection.
